\chapter{Library Ecosystem and Tooling}
\label{chap:ch6}

\par Modern backend applications often expose the same HTTP and JSON interface, even if they are implemented in very different languages. In Gentlix Bank, the C and Rust backends look identical from the frontend point of view, but the internal libraries, tools, and amount of custom infrastructure code are very different.

\section{The Database Layer}
\label{sec:ch6sec1}

\par Both backends use PostgreSQL as the database, but the abstractions around it are not comparable. In the C version the only external library is \texttt{libpq}, the official PostgreSQL C client, which exposes a raw connection handle and string-based query functions. There is no ORM, no query builder and no type-safe mapping, so almost all structure above the wire protocol was implemented manually.

\par In contrast, the Rust backend is built on the \texttt{sqlx} crate, which provides an asynchronous connection pool, compile-time checked SQL queries and automatic mapping between rows and Rust structs. The database layer in Rust is still explicit and under control, but many pieces that were hand-written in C appear here as features of the standard ecosystem.

\subsection{Low-level driver and manual \texorpdfstring{\texttt{Db}}{Db} wrapper in C}
\label{subsec:ch6sec1-1}

\par In the C implementation, database access starts directly from \texttt{libpq}. Each SQL statement is a raw string with positional placeholders, and every parameter is passed as a \texttt{const char *}. The library returns an opaque \texttt{PGresult} object, from which columns are read one by one with \texttt{PQgetvalue(res, row, col++)} and converted manually to C types. All resource management, including connection lifetime and \texttt{PQclear} on every exit path, is the developer's responsibility.

\par \textbf{Schema alignment.} Does column five still mean \texttt{balance\_cents}? In C 
the answer is a correct \texttt{PQgetvalue} index. \texttt{account\_repo.c} maps rows with 
a running \texttt{col} counter:

\begin{lstlisting}[language=C, caption={account\_repo.c --- Column order must match SELECT
list order; a migration that reorders columns compiles but mis-loads balances.}]
const char *id_s           = PQgetvalue(res, row, col++);
const char *user_id_s      = PQgetvalue(res, row, col++);
const char *account_type_s = PQgetvalue(res, row, col++);
const char *currency_s     = PQgetvalue(res, row, col++);
const char *balance_s      = PQgetvalue(res, row, col++);
const char *iban_s         = PQgetvalue(res, row, col++);
\end{lstlisting}

\par Strings from PostgreSQL pass through \texttt{util\_str\_to\_i64} and domain parsers 
(\texttt{account\_type\_from\_str}, \texttt{iban\_parse}). A \texttt{NULL} field from 
\texttt{PQgetvalue} becomes \texttt{0} if the code forgets \texttt{PQgetisnull}---accounts 
can appear with a zero balance while HTTP still returns 200. Gatchell~\cite{Gatchell2003} 
stresses that C database code is really string handling plus discipline; Gentlix adds 
\texttt{RepoError} enums so failures at least map to structured service errors when 
developers remember to check return values.

\par \textbf{Ledger atomicity.} Money-moving operations call 
\texttt{begin}/\texttt{commit}/\texttt{rollback} on the repository vtable so debit, credit, 
and insert share one PostgreSQL transaction. The C service frees in-memory rows and rolls 
back on every failure path; Rust uses \texttt{?} and \texttt{Drop} for the same guarantee 
(Section~\ref{sec:ch5sec2}). That behaviour is custom glue in \texttt{db.c}, not a libpq 
feature, but it is essential so a mid-write crash cannot leave a half-updated ledger.

\par To make libpq usable day to day, the project introduces a custom \texttt{Db} abstraction. 
This wrapper hides the \texttt{PGconn} pointer, provides helpers to connect with a URL and 
execute parametrized statements, and normalises errors into \texttt{RepoError}. This 
infrastructure is not part of the business domain; it exists because the C driver is too 
raw for higher layers.

\subsection{Type-checked queries and \texorpdfstring{\texttt{sqlx}}{sqlx} in Rust}
\label{subsec:ch6sec1-2}

\par The Rust backend uses \texttt{sqlx} with the PostgreSQL feature and a \texttt{PgPool} 
connection pool. Queries use macros such as \texttt{sqlx::query!} or \texttt{sqlx::query\_as!}, 
verified against the live schema during development or against a cached description in the 
\texttt{.sqlx} directory for offline CI builds.

\par Palmieri~\cite{Palmieri2022} recommends compile-time SQL checking so schema drift fails 
the build instead of corrupting production rows. Gentlix follows that approach: repository 
functions map rows directly into domain structs; if a migration renames or reorders a column, 
the macro no longer compiles until the query is fixed. Drysdale~\cite{Drysdale2024} describes 
the same idea as pushing validation left---errors appear where developers already look 
(the compiler) rather than in a customer's statement view.

\par \textbf{Typed rows versus manual indices.} Where C increments \texttt{col++}, Rust 
binds names in the \texttt{SELECT} list to struct fields. Optional database columns become 
\texttt{Option<T>} and force explicit handling---the Rust equivalent of remembering 
\texttt{PQgetisnull}, but enforced on every query using the macro.

\par When a query is valid, \texttt{sqlx} maps each row into a strongly typed struct. 
Connection management is handled by \texttt{PgPool}, which clones cheaply and shares safely 
across async tasks without \texttt{unsafe} in application code. Blackburn~\cite{Blackburn2025} 
groups ``wrong column index'' with type-layer failures; \texttt{sqlx} moves that class of bug 
from runtime review into compile time for the Rust backend.

\par The trade-off is tooling dependency: developers need PostgreSQL (or \texttt{.sqlx} cache) 
at build time. C only needs headers and link libraries, but pays with 29 manual 
\texttt{PQgetvalue} sites counted in the Phase~1 audit (Chapter~\ref{subsec:ch8sec1audit}).

\section{The HTTP Server Layer}
\label{sec:ch6sec2}

\par Both backends expose the same REST endpoints, accept and return JSON, and implement the same authentication rules, but library support differs sharply. The C backend sits on libmicrohttpd with a custom router; the Rust backend uses \texttt{axum} with typed extractors and \texttt{serde}.

\subsection{Embedded C HTTP engine and custom router}
\label{subsec:ch6sec2-1}

\par In C, libmicrohttpd opens sockets, parses HTTP/1.1, and delivers each request to a callback with raw buffers. The project adds a routing layer for paths and methods, plus jansson for JSON field-by-field parsing. Even a simple authenticated route passes through several hand-written layers before reaching a service function.

\subsection{Rust \texorpdfstring{\texttt{axum}}{axum} framework and extractors}
\label{subsec:ch6sec2-2}

\par Rust routing is declarative. \texttt{router.rs} nests resource modules under \texttt{/api}, attaches JWT middleware, and maps paths to handler functions whose signatures the compiler checks:

\begin{lstlisting}[language=Rust, caption={router.rs --- The HTTP surface mirrors the C
backend; only one backend runs at a time on port 6767.}]
Router::new()
    .route("/health", get(health))
    .nest("/api/users",        routes::users::router(state.clone()))
    .nest("/api/accounts",     routes::accounts::router(state.clone()))
    .nest("/api/affiliates",   routes::affiliates::router(state.clone()))
    .nest("/api/transactions", routes::transactions::router(state.clone()))
    .nest("/api/dashboard",    routes::dashboard::router(state.clone()))
\end{lstlisting}

\par Request bodies deserialize into structs with \texttt{\#[derive(Deserialize)]}; responses use \texttt{Serialize}. Shared \texttt{AppState} holds pools and services instead of threading pointers through every function.

\subsection{Custom infrastructure: JWT and domain utilities}
\label{subsec:ch6sec2-3}

\par JWT handling illustrates ecosystem gaps in C. The C backend builds headers and payloads, base64url-encodes them, and signs with OpenSSL HMAC primitives in roughly 294 lines of custom code. Rust delegates signing and verification to \texttt{jsonwebtoken} after defining a claims struct. IBAN validation follows the same pattern: manual MOD-97 code in C versus a validated \texttt{IBAN} type in Rust (Section~\ref{sec:ch5sec1}).

\subsection{Error responses at the HTTP boundary}
\label{subsec:ch6sec2-5}

\par Service errors stop at \texttt{ServiceError}; the HTTP layer converts them to status 
codes and JSON the React app already handles. Both backends target the same envelope 
(\texttt{status}, \texttt{code}, \texttt{message}). Insufficient funds become HTTP~400 
\texttt{validation\_error}; forbidden transfers become 403; not-found accounts become 404.

\par C centralises mapping in \texttt{http\_error.c}; handlers in \texttt{http\_transactions.c} 
check the service \texttt{bool} and call \texttt{http\_error\_from\_service\_error}:

\begin{lstlisting}[language=C, caption={http\_error.c --- HTTP mapping lives only here,
not in the service layer (Section~\ref{sec:ch5sec2}).}]
case SERVICE_ERROR_VALIDATION:
    body.status = 400;
    body.code   = "validation_error";
    break;
\end{lstlisting}

\par Rust wraps \texttt{ServiceError} in \texttt{ApiError} and implements 
\texttt{IntoResponse} in \texttt{server/error.rs}. Handlers use \texttt{?}; axum turns 
the error into JSON automatically. Palmieri~\cite{Palmieri2022} stresses that this 
boundary is where unhandled failures become user-visible incidents---especially when 
HTTP~200 would mask wrong financial state.

\par Request parsing differs too: C reads query strings with \texttt{MHD\_lookup\_connection\_value} 
and \texttt{strtoll}; Rust deserialises into typed DTOs (Section~\ref{sec:ch6sec4}). 
The contract matches; only the proof burden differs.

\subsection{Concurrency model and async runtime}
\label{subsec:ch6sec2-4}

\par The C server uses blocking I/O on libmicrohttpd worker threads. The Rust server runs 
on \texttt{tokio}; handlers \texttt{await} \texttt{sqlx} calls. CPU-heavy work---Argon2 
login, statement replay, analytics aggregation---runs inside \texttt{spawn\_blocking} in 
Rust but on the same thread that accepted the request in C. Flitton and Morton~\cite{Flitton2023} 
explain why blocking work must leave the async executor; the mutex/atomic/thread design 
inside analytics is described once in Section~\ref{sec:ch5sec3}.

\section{The REST Contract and Pagination}
\label{sec:ch6sec4}

\par Frontend parity depends on more than similar page layouts. Both backends must agree on paths, query parameter names, JSON field names, and error shapes. Gentlix exposes 22 authenticated routes under \texttt{/api} plus \texttt{GET /health}. Table~\ref{tab:api_routes} groups them by resource; method and path combinations match in C and Rust.

\begin{table}[htbp]
\centering
\small
\begin{tabular}{|l|l|p{6.5cm}|}
\hline
\textbf{Group} & \textbf{Prefix} & \textbf{Main operations} \\
\hline
Users & \texttt{/api/users} &
Register, login, profile with accounts \\
\hline
Accounts & \texttt{/api/accounts} &
Open account, check availability, get by id \\
\hline
Affiliates & \texttt{/api/affiliates} &
List/create, get/update/delete, resolve transfer target \\
\hline
Transactions & \texttt{/api/transactions} &
Deposit, withdrawal, send, transfer, payment; recent list; \\
& & statement; summary; monthly summary \\
\hline
Dashboard & \texttt{/api/dashboard} &
Aggregated cards, recent activity, chart data \\
\hline
Health & \texttt{/health} &
Unauthenticated liveness probe \\
\hline
\end{tabular}
\caption{REST surface shared by both Gentlix backends (22 \texttt{/api} routes + health)}
\label{tab:api_routes}
\end{table}

\par \textbf{Pagination.} List endpoints that can return long histories accept 
\texttt{limit} and \texttt{offset} query parameters. The account statement is the 
richest example: the service replays \emph{all} matching transactions to compute 
correct running balances, then returns only the requested page together with metadata 
the UI needs to draw pager controls without guessing.

\par Query parameters on \texttt{GET /api/transactions/statement} include:
\begin{itemize}
  \item \texttt{account\_id} --- required account scope;
  \item \texttt{from}, \texttt{to} --- optional date bounds (\texttt{YYYY-MM-DD});
  \item \texttt{limit}, \texttt{offset} --- page size and skip (defaults applied in the service layer);
  \item \texttt{transaction\_type} --- optional filter (deposit, withdrawal, send, etc.);
  \item \texttt{sort} --- \texttt{oldest} or \texttt{newest}.
\end{itemize}

\par The JSON response includes \texttt{items} (each with \texttt{balance\_after\_cents}), 
\texttt{total\_count}, and optional \texttt{opening\_balance\_cents} / 
\texttt{closing\_balance\_cents} for the filtered window. Recent transactions on 
\texttt{GET /api/transactions/recent} use the same \texttt{limit}/\texttt{offset} pair 
with a smaller default page size.

\par Monthly analytics (\texttt{GET /api/transactions/summary/monthly}) uses 
\texttt{per\_account\_limit} to cap how many rows are loaded per owned account before 
parallel aggregation---the k6 W3 workload sets this high to stress CPU and memory.

\par C parses query strings manually from the libmicrohttpd connection; Rust deserialises 
into a typed DTO. The contract is identical; the proof burden is not:

\begin{lstlisting}[language=C, caption={http\_transactions.c --- Statement pagination
parameters are read as strings and converted to integers.}]
const char *limit_s  = MHD_lookup_connection_value(
    conn, MHD_GET_ARGUMENT_KIND, "limit");
const char *offset_s = MHD_lookup_connection_value(
    conn, MHD_GET_ARGUMENT_KIND, "offset");
int64_t limit  = limit_s  ? strtoll(limit_s,  NULL, 10) : 0;
int64_t offset = offset_s ? strtoll(offset_s, NULL, 10) : 0;
\end{lstlisting}

\begin{lstlisting}[language=Rust, caption={transactions.rs (routes) --- The same
parameters are validated when the handler is type-checked.}]
pub struct AccountStatementQueryDto {
    pub account_id: i64,
    pub from: Option<String>,
    pub to: Option<String>,
    pub limit: Option<i64>,
    pub offset: Option<i64>,
    pub transaction_type: Option<String>,
    pub sort: Option<String>,
}
\end{lstlisting}

\section{Build and Dependency Management}
\label{sec:ch6sec3}

\par The C backend is built with CMake, which locates system libraries (\texttt{libpq}, libmicrohttpd, jansson, OpenSSL, Argon2) and links them explicitly. The Rust backend uses \texttt{cargo} with dependencies declared in \texttt{Cargo.toml}. Chapter~\ref{sec:ch8sec1} compares build reproducibility, manual audit counts, and undefined-behaviour risk in more detail; here the emphasis is on day-to-day workflow---one \texttt{cargo build} versus installing MSYS2 packages and tuning \texttt{CMakeLists.txt} on each machine.

\section*{Summary}

\par From the frontend perspective, the C and Rust implementations of Gentlix Bank are 
indistinguishable. Inside, libpq indices versus \texttt{sqlx} macros, custom JWT versus 
crates, and manual versus typed HTTP parsing differ---but each topic is treated in one 
place: database shape in Section~\ref{sec:ch6sec1}, HTTP errors and runtime in 
Section~\ref{sec:ch6sec2}, pagination in Section~\ref{sec:ch6sec4}, domain invariants 
in Section~\ref{sec:ch5sec1}.
